\documentclass{report}
\usepackage{amsmath,amssymb}
\setlength{\parindent}{0mm}
\setlength{\parskip}{1em}
\begin{document}
\begin{center}
\rule{6in}{1pt} \
{\large Sou-Cheng Choi \\
{\bf CS-MINRES: A Krylov Subspace Method for Complex Symmetric Linear Systems and Least-Squares Problems}}

Computation Institute \\ University of Chicago \\ Searle Chemistry Laboratory \\ 5735 South Ellis Avenue \\ Chicago \\ IL 60637
\\
{\tt sctchoi@uchicago.edu}\end{center}

While there is no lack of high-quality Krylov subspace solvers for
(complex) Hermitian systems, there are few for complex symmetric systems,
which have become increasingly important in modern applications including
quantum dynamics, electromagnetics, power systems, and many more.

For a large consistent complex symmetric system $Ax=b$, i.e., $A=A^T$,
one may apply a non-Hermitian Krylov subspace method such as BiCG, CGS,
or GMRES, overlooking the symmetry of $A$, or a Hermitian Krylov solver
such as CG, SYMMLQ, or MINRES on the equivalent normal equation or an
augmented system twice the original dimension. These have the
disadvantages of increasing either memory, conditioning, or computational
costs. An exception is a special version of QMR by Freund (1992), but
that may be affected by non-benign breakdowns unless look-ahead is
implemented; furthermore, it is designed for consistent problems only.

We describe CS-MINRES, a new algorithm for solving linear systems and
least-squares problems with complex symmetric coefficient matrices. In
all cases, CS-MINRES computes the unique minimum-length, i.e.,
pseudoinverse, solution. It may be regarded as an extension of MINRES by
Paige and Saunders (1975) or the more recent MINRES-QLP by Choi, Paige
and Saunders (2011). CS-MINRES does not suffer from non-benign breakdowns
and is immediately applicable to real symmetric problems, in which case
it is mathematically equivalent to MINRES-QLP.

Like MINRES and MINRES-QLP, CS-MINRES has elegant theoretical properties,
and is sufficiently robust for ill-conditioned and ill-posed problems. We
will present extensive numerical experiments to demonstrate the
scalability and robustness of CS-MINRES.


\end{document}
